% =====================================================================
% Paste-ready captions for Figures 4-6, revised against the current run.
%
% Generated alongside scripts/plots/. Every number below was computed from
% the rebuilt community_results/ aggregate (build_community.py, 2026-08-31
% 18:16). Regenerate if the runs change:
%   cd scripts/plots && python run_all.py
%
% ALL Load Duration Curve numbers here are on the "mean_day" basis: the
% representative daily curve (96 fifteen-minute bins) sorted descending.
% That is what validate_simulation.py's plot_load_duration_curve and
% calculate_ldc_metrics use, and what validate_community.py uses via
% representative_daily_curve(), so the figures match the draft
% val_ldc_*.png panels and the metrics tables.
%
% Do NOT quote these alongside a --ldc-basis pooled figure. Pooling every
% sample instead of averaging to a daily curve gives a much wider curve:
% P1 measured reaches 7.50 W with a P95 of 4.66 W, against 4.90 W and
% 4.22 W on the mean_day basis.
%
% Figure files produced by run_all.py:
%   figure4_profile_load_curves.png
%   figure5_profile_ldc.png
%   figure6_community_ldc_envelope.png
%
% WHAT TO REMOVE FROM main.tex
% ----------------------------
%  1. fig:tier1_overlay  — drop the \fbox placeholder, point
%     \includegraphics at figure4_profile_load_curves.png, and add "for
%     May" to the caption (the prose at l.839 already fixes the window to
%     May; the caption currently omits it).
%  2. fig:ldc — drop the four \begin{subfigure} blocks pointing at
%     validate_simulation.py's val_ldc_*.png. figure5_profile_ldc.png is a
%     single 2x2 figure and replaces all four. Delete the "Provisional
%     styling — to be replaced with journal-formatted version" note; these
%     panels are journal-styled and generated against this run.
%  3. fig:surveybased_vs_sociotechnical — this is a Profile-2 three-way
%     comparison (Measured vs Survey-Based vs Socio-Technical). The
%     community-aggregate Figure 6 does not exist in main.tex yet, so it
%     needs a new figure environment; decide whether Fig. 6 replaces the
%     Profile-2 figure or sits alongside it.
%  4. fig:tier2_heatmap is still a placeholder and is not in the Fig. 4-6
%     set. Either generate it or cut it — it currently shifts the figure
%     numbering relative to the numbering used here.
%
% WINDOW WARNING — applies to Figure 6
% ------------------------------------
% Figures 4 and 5 are May panels. Figure 6 defaults to the FULL YEAR, and
% the strength of its contrast depends heavily on that choice
% (evening-peak error is +13.6% over the year but -4.4% in May, i.e. it
% changes sign). Whichever you choose, state the window in the caption —
% the captions below do. Do not silently mix windows across the three.
% =====================================================================


% ---------------------------------------------------------------- Fig. 4
\caption{Measured versus simulated (Socio-Technical Model) mean daily load
curves for May, by Energy Behavior Profile. Measured telemetry is shown as
a solid black line, the Socio-Technical Model as a dashed navy line. May is
the month with continuous telemetry across the logger set, and falls between
the community's two principal mobility peaks (growing season, February--April;
free-grazing and migration, July--September), which minimises the confounding
effect of household absence. Panels are not equally weighted: P1, P3 and P4
each pool two instrumented households and P2 a single one, and each panel
states its own sample size and number of measured days. Note that P2 rests on
30 rather than 31 days, its telemetry being absent for 31 May and most of
30 May.}
\label{fig:tier1_overlay}


% ---------------------------------------------------------------- Fig. 5
\caption{Load Duration Curve by Energy Behavior Profile for May: measured
(solid black) and Socio-Technical simulated (dashed navy) power, each sorted
in descending order against the fraction of the day spent at or above that
level. Discarding timing isolates how well the model reproduces the
distribution of demand levels a household occupies, independent of when each
level occurs --- the quantity that governs inverter and battery sizing. Each
series is the representative daily curve (96 fifteen-minute bins, averaged
across the days of the window) sorted descending, the same construction as
the LDC-RMSE and P95 values in Table~\ref{tab:ldc_summary}, so figure and
table describe the same object. Annotations give the P95 level, the demand
exceeded only 5\% of the day.}
\label{fig:ldc}


% ---------------------------------------------------------------- Fig. 6
\caption{Community aggregate over the full simulated year: Socio-Technical
(heterogeneous, solid blue) reference versus Survey-Based (homogeneous, dashed
green) Model. (a) Mean daily load curve for each community (heavy lines) drawn
over its individual daily traces (thin lines) and its 5th--95th percentile
envelope across days. Both models place the evening peak of the mean curve at
20:00, so peak \emph{timing} is not what separates them, though the
Survey-Based Model overstates its height by 27.9\% (229.5~W against 293.6~W);
what separates them is how each distributes demand across the rest of the day.
The Survey-Based Model understates the overnight base load by 48.5\% (77.5~W
against 39.9~W for 00:00--04:00), overstates the morning window by 80.0\%
(54.9~W against 98.9~W for 04:00--10:00) and the daytime window by 54.8\%
(44.6~W against 69.0~W for 10:00--17:00), and roughly doubles the morning peak
of the mean curve (84.0~W at 05:45 against 169.7~W at 06:00). Because it is
built from a single static year-round average, its curve is nearly identical
in every month, whereas the heterogeneous reference varies seasonally; the
overstatement of aggregate mean power is therefore 22.1\% across the year but
only 3.1\% within May alone. The two communities' day-to-day \emph{spread}, by
contrast, is comparable in magnitude (mean envelope width 51.1~W heterogeneous
against 38.5~W homogeneous over the year, and 34.2~W against 35.5~W within
May); what distinguishes them is the tail. Correlating each day against the
mean normalised daily shape, the most atypical heterogeneous day reaches 0.911
whereas the homogeneous Model never falls below 0.959: the heterogeneous
community generates genuinely aberrant days, and the homogeneous one does not.
(b) Load Duration Curve, each series sorted descending against the fraction of
the day at or above that level: the homogeneous Model sits above the reference
across most of the distribution, with a P95 of 259.7~W against 204.6~W and a
maximum of 293.6~W against 229.5~W.}
\label{fig:community_ldc_envelope}


% =====================================================================
% CLAIMS THAT DID NOT SURVIVE THE CURRENT RUN — do not reinstate
% =====================================================================
% "the heterogeneous reference spreads across a wide day-to-day band, the
%  homogeneous Model a narrow, near-deterministic one"
%
% Not supported. Mean p5-p95 envelope widths, rebuilt heterogeneous vs
% homogeneous: 51.1 vs 38.5 W over the full year (1.33x) and 34.2 vs
% 35.5 W within May (0.96x, i.e. reversed). Level-normalised shape
% coherence is likewise close: mean correlation to the mean daily shape
% 0.975 vs 0.985 over the year, 0.984 vs 0.986 in May.
%
% The rebuild does not rescue the claim: the stale aggregate gives 52.6 W
% and 1.37x, essentially the same answer.
%
% The likely origin of the error is MRSD (0.150 heterogeneous vs 0.016
% homogeneous, Table \ref{tab:...}). MRSD measures WITHIN-day roughness,
% not across-day spread. The homogeneous Model is indeed smooth within a
% day, but RAMP redraws switch-on times for every household on every day,
% so its day-to-day envelope does not collapse. MRSD cannot be cited in
% support of a day-to-day variability claim.
%
% Fig. 6(a) above therefore claims the tail (min shape correlation), which
% the data does support, rather than the band width, which it does not.
